\subsection{The Thom Isomorphism}

Recall that a vector bundle of rank $n$ is a map $\pi\colon E\to M$ such that the fibers $E_x=\pi^{-1}x$ all have the structure of a vector space.
Furthermore there is an open cover $\set{U_i}$ of $M$ and local trivializations $\phi_i\colon E|_{U_i}=\pi^{-1}U_i\isoto U_i\times\bR^n$ which is linear
on the fibers (meaning $E_x$ is mapped linearly to $\set x\times\bR^n$).
Then the maps
$$ \phi_i\circ\phi_j^{-1}\colon(U_i\cap U_j)\times\bR^n \isoto (U_i\cap U_j)\times\bR^n $$
are linear isomorphisms on the fibers and thus of the form
$$ \phi_i\circ\phi_j^{-1}(p,v) = (p,g_{ij}(p)v) $$
for {\emphcolor transition functions}
$$ g_{ij} \colon U_i\cap U_j \to \gl_n(\bR) . $$
The transition functions satisfy the {\emphcolor cocycle condition}:
$$ g_{ij}\circ g_{jk} = g_{jk}\qquad\text{on $U_i\cap U_j\cap U_k$.} $$

If $g'_{ij}$ are other transition functions coming from another trivialization $\phi'_i$, then notice that $\phi_i\circ(\phi'_i)^{-1}$ is of the form
$(p,v)\mapsto(p,\lambda_i(p)v)$ for $\lambda_i\colon U_i\to\gl_n(\bR)$.
Thus
$$ g_{ij} = \phi_i\phi_j^{-1} = \lambda_i\phi'_i(\phi'_j)^{-1}\lambda_j^{-1} = \lambda_ig'_{ij}\lambda_j^{-1} . $$
And so we see that all transitions are equivalent in this sense.

\bdefn

    Suppose $E$ is a vector bundle which has a trivialization whose transition functions take values in a subgroup $H\leq\gl_n(\bR)$.
    Then we say that {\emphcolor $E$'s structure group may be reduced to $H$}.

    The vector bundle $E$ is said to be {\emphcolor oriented} by trivializations $\phi_i$ if $g_{ij}$ has positive determinant.
    That is to say, its structure group may be reduced to $\gl_n^+(\bR)$.
    If $\psi_j$ is another trivialization, it is said to be {\emphcolor equivalent} to $\phi_i$ if $\phi_i\circ\phi_j^{-1}(x)\in\gl_n(\bR)$
    has positive determinant.

\edefn

We easily see that $M$ is orientable as a manifold iff its tangent bundle $TM$ is orientable as a vector bundle.

\blemm

    The structure group of any vector bundle $E\to M$ may be reduced to $\orth(n)$.
    It may be reduced to $\sorth(n)$ iff $E$ is orientable.

\elemm

\bproof

    Give $E$ a Riemannian metric structure (which exists by gluing via a partition of unity).
    Relative to this metric, cover $E$ with orthonormal trivializations (which exist by the Gram-Schmidt process).
    Then $g_{ij}$ preserves orthonormal bases, so it takes values in $\orth(n)$.
    \qed

\eproof

Recall that if $f\colon M\to N$ is a smooth map and $\pi\colon E\to N$ is a vector bundle, we may pull back to form a vector bundle $f^*\pi\colon f^*E\to M$.
Explicitly set
$$ f^*E = \set{(p,v)\in M\times E}[f(p)=\pi(v)] = \coprod_{p\in M}E_{f(p)},\qquad f^*\pi(p,v) = p . $$
Its fibers are $E_{fp}$.
If $\phi_i\colon E|_{V_i}\to V_i\times\bR^n$ is a trivialization of $E$, then define $U_i=f^{-1}V_i$ and
$$ \psi_i\colon(f^*E)|_{U_i} = \set{(p,v)\in f^*E}[f(p)\in V_i] \isoto U_i\times\bR^n,\qquad \psi_i(p,v) = (p,\phi_i(fp,v)) . $$
The transition functions of the pullback are $f^*g_{ij}=g_{ij}\circ f$.

Suppose $E\to M$ and $F\to M$ are vector bundles.
Then we may construct the vector bundle $\hom(E,F)\to M$ whose fibers are $\hom(E_p,F_p)$.
Define $\iso(E,F)\to M$ to be the fiber bundle whose fibers are isomorphisms $E_p\isoto F_p$.
This is a fiber bundle with fiber $\gl_n(\bR)$.
Furthermore an isomorphism of vector bundles $E\isoto F$ is simply a section of $\iso(E,F)$.

\bprop[name={bundlehomotopy}]

    Suppose $f_0,f_1\colon M\to N$ are homotopic, and $E\to N$ is a vector bundle.
    Then $f_0^*E$ and $f_1^*E$ are isomorphic vector bundles.

\eprop

%\bnote
%
%    A paracompact space is one for which every open cover has a locally finite refinement.
%    Locally finite means that every point in $M$ has a neighborhood which meets only finitely many opens in the refinement.
%    Every metric space, and thus every manifold, is paracompact, so the assumption here is superfluous.
%
%\enote

\bproof

    Let $f\colon M\times[0,1]\to N$ be a homotopy from $f_0$ to $f_1$, and let $\pi\colon M\times[0,1]\to M$ be the projection.
    Suppose that for some $t_0\in[0,1]$, $f_{t_0}^*E$ is isomorphic to a vector bundle $F\to M$.
    We will then show that for all $t$ sufficiently close to $t_0$, $f_t^*E$ is isomorphic to $F$.
    By compactness, we get that {\it all\/} $f_t^*E$ are isomorphic to $F$; taking $F=f_1^*E$ we complete the proof.

    Over $M\times[0,1]$ we have the bundles $f^*E$ and $\pi^*F$.
    Because $f_{t_0}^*E\cong F$, we have a section of $\iso(f^*E,\pi^*F)$ over $M\times\set{t_0}$.
    This is a section of $\hom(f^*E,\pi^*F)$, which is a vector bundle, and thus this section may be extended to a neighborhood of $M\times\set{t_0}$.
    Furthermore it must be invertible in a neighborhood of $M\times\set{t_0}$.
    So this extends to a local section of $\iso(f^*E,\pi^*F)$ around $M\times\set{t_0}$.
    This gives us precisely what we want.
    \qed

\eproof

\bcoro

    If $M$ is contractible, all of its vector bundles are trivial.

\ecoro

\bproof

    Let $f\colon M\to*\inj M$ be a constant map, so that $\id\colon M\to M$ is homotopic to it.
    Then for a vector bundle $E\to M$, $E$ and $f^*E$ are isomorphic and $f^*E$ is just the trivial bundle $M\times E_*$.
    \qed

\eproof

\blemm[name={orientedvboveroriented}]

    If $E\to M$ is an oriented vector bundle such that $M$ is also orientable, then $E$ is orientable as a manifold.
    Furthermore if $M$ is of finite type, so too is $E$.

\elemm

\bproof

    Let $(U_i,\psi_i)$ be an oriented atlas for $M$, so that $h_{ij}=\psi_i\circ\psi_j^{-1}$ has Jacobian with positive determinant.
    And let $\phi_i\colon E|_{U_i}\to U_i\times\bR^n$ be local trivializations of $E$ such that $g_{ij}$ have positive determinant.
    Now
    $$ E|_{U_i} \xtos{\phi_i} U_i\times\bR^n \xtos{\psi_i\times\id} \bR^n\times\bR^n $$
    gives an atlas for $E$.
    The transition functions of this atlas are
    $$ (\psi_i\times\id)\circ\phi_i\circ\phi_j^{-1}\circ(\psi_j^{-1}\times\id) \colon \bR^{n+m} \to \bR^{n+m} , $$
    which send $(x,y)$ to $(h_{ij}(x),g_{ij}(\psi_i(x))y)$.
    This has Jacobian
    $$ \pmatrix{J(h_{ij})&*\cr0&g_{ij}\circ\psi_i} , $$
    which has positive determinant, as desired.

    The second claim is clear, as $E|_{U_i}\cap E|_{U_j}=E|_{U_i\cap U_j}\cong(U_i\cap U_j)\times\bR^n$, so that we see $E|_{U_i}$ forms
    a good cover if $U_i$ does.
    \qed

\eproof

\bcoro

    Suppose $M$ is an orientable $m$-manifold and $E\to M$ is an orientable rank-$n$ vector bundle.
    Suppose further that $M$ is of finite type.
    Then
    $$ H^\bul_\dR(E) \cong H^\bul_\dR(M),\qquad H^\bul_c(E) \cong H^{\bul-n}(M) . $$

\ecoro

\bproof

    The first isomorphism comes from the fact that the zero section $M\to E$ embeds $M$ into $E$, and forms a deformation
    retract of $E$ onto $M$ (by connecting a vector $(p,v)\in E$ to $(p,0)$ along a line).
    Thus $E$ is homotopy equivalent to $M$.

    For the second isomorphism we have by the above lemma that $E$ is orientable.
    Then by Poincar\'e duality and the first isomorphism
    $$ H^\bul_c(E) \cong (H^{n+m-\bul}_\dR(E))^* \cong (H^{n+m-\bul}_\dR(M))^* \cong H^{\bul-n}_c(M) . \qed $$

\eproof

\bdefn

    Let $E\to M$ be a vector bundle.
    A differential form $\omega\in\Omega^\bul(E)$ has {\emphcolor compact support vertically} if for every compact $K\subseteq M$,
    $\pi^{-1}K\cap\supp\omega\subseteq E$ is compact.
    Note then that $\supp\omega|_{E_x}\subseteq E_x\cap\supp\omega$ is closed in a compact set, thus compact.
    So if $\omega$ has compact support vertically, while it may not have compact support, it does have compact support in $E$'s fibers.

    Denote by $\Omega_\cv^\bul(E)$ the space of differential forms on $E$ with compact vertical support.

\edefn

Let $E=M\times\bR^n$ be a trivial bundle, and denote by $t_1,\dots,t_n$ the coordinates of $\bR^n$.
Then if we let $x_1,\dots,x_m$ be local coordinates of $M$, we have that a differential form on $E$ looks locally like
$$ \omega = \sum_{I,J}\alpha_{I,J}\d x_I\wedge\d t_J,\qquad \alpha_{I,J} \in C^\infty(M\times\bR^n) . $$
Then
$$ \omega = \sum_{I,J}\alpha_{I,J}(\pi^*\d x_I)\wedge\d t_J , $$
were $\pi\colon E\to M$ is the projection.
Thus differential forms on $E$ are linear combinations of forms of the form $(\pi^*\tau)\wedge f(x,t)\d t_J$ where $\tau$ is a differential
form on $M$ and $f(x,t)\in C^\infty(E)$.

If $\omega$ has vertical compact support then for each $x\in M$, $f(x,t)$ has compact support in $\bR^n$.
Thus we may define the map $\pi_*\colon\Omega^\bul_\cv(E)\to\Omega^{\bul-n}(M)$ by considering these forms.
\benum
    \item If $\omega=(\pi^*\tau)\wedge f(x,t)\d t_J$ where $J$ is not of length $n$, then $\pi_*\omega=0$.
    \item If $\omega=(\pi^*\tau)\wedge f(x,t)\d t_1\wedge\cdots\wedge\d t_n$, then
    $$ \pi_*\omega = \parens{\int_{\bR^n}f(x,t)\d t_1\cdots\d t_n}\tau . $$
\eenum

In general if $E\to M$ is an oriented vector bundle, let $(U_i,\phi_i)$ be an oriented trivialization then let $\bar x,\bar y$ be coordinates
for $U_i,U_j$ (in $M$) and $\bar t,\bar u$ be frames for $E|_{U_i}$ and $E|_{U_j}$ given by these trivializations.
Because the trivialization is oriented, for every point of $U_i\cap U_j$, $\bar t$ and $\bar u$ are related by an element of $\gl_n^+(\bR)$.
A form on $E$ is locally of the form $\omega|_{U_i}=(\pi^*\tau)\wedge f(x,t)\d t_J$.
Again if $J$ does not contain all $1,\dots,n$ then we define $\pi_*\omega|_{U_i}=0$.
And otherwise $\omega_i=\omega|_{U_i}=(\pi^*\tau)f(x,t)\d t_1\wedge\cdots\wedge\d t_n$ so define
$$ \pi_*\omega_i = \parens{\int_{\bR^n}f(x,t)\d t_1\cdots\d t_n}\tau . $$
Suppose that in $U_j$, $\omega_j=(\pi^*\eta)g(x,u)\d u_1\wedge\cdots\wedge\d u_n$ and so
$$ \pi_*\omega_j = \parens{\int_{\bR^n}g(x,t)\d u_1\cdots\d u_n}\eta . $$
Let $\Phi\colon U_i\cap U_j\to\gl_n^+(\bR)$ be the transition from $\bar u$ to $\bar t$ so that $\bar u=\Phi(x)\bar t$.
We have that in $U_i\cap U_j$,
$$ \omega_j = (\pi^*\eta)g(x,u)\d u_1\wedge\cdots\wedge\d u_n = (\pi^*\eta)g(x,\Phi(x)t)\det(\Phi(x))\d t_1\wedge\cdots\wedge\d t_n
= (\pi^*\tau)f(x,t)\d t_1\wedge\cdots\wedge\d t_n = \omega_i . $$
Integrating over $E_x$ for a set $x$ with respect to $\bar t$ gives us
$$ \tau\cdot\int_{\bR^n}f(x,t)\d\bar t = \eta\cdot\int_{\bR^n}g(x,\Phi(x)t)\det{\Phi(x)}\d\bar t . $$
Because $\Phi(x)$ has positive determinant this is equal to
$$ = \eta\cdot\int_{\bR^n}g(x,u)\d\bar u . $$
Thus the $\pi_*\omega_i$ all agree on their intersections and glue together to give $\pi_*\omega\in\Omega^{\bul-n}(M)$.

\bdefn

    Let $E\to M$ be an oriented vector bundle.
    The map defined above
    $$ \pi_*\colon\Omega^\bul_\cv(E)\to\Omega^{\bul-n}(M) $$
    is called {\emphcolor integration along the fiber}.

\edefn

\bprop

    Integration along the fiber is a map of complexes.
    Thus it descends into a map of cohomology
    $$ \pi_*\colon H^\bul_\cv(E) \to H^{\bul-n}_\dR(M) . $$

\eprop

\bproof

    We must show that $\pi_*$ commutes with exterior differentiation.
    Let $(U_i,\phi_i)$ be a trivialization for $E$, with a partition of unity $\rho_i$ subordinate to it.
    Then for $\omega\in\Omega^\bul_\cv(E)$, we have $\omega=\sum_i\rho_i\omega$.
    Because both $\pi_*$ and $\d$ are linear, it is sufficient to prove this for $\rho_i\omega$.
    Thus we may assume from the outset that $E=M\times\bR^n$.

    If $\omega=(\pi^*\tau)f(x,t)\d t$, then
    $$ \d\pi_*\omega = \d\parens{\int f(x,t)\d t} = \parens{\sum_i\d x_i\int\pdvof f{x_i}(x,t)\d t}\tau
    + (\d\tau)\int f(x,t)\d t . $$
    And
    $$ \multline{
        \pi_*\d\omega = \pi_*((\pi^*\d\tau)f\d t+\d f\wedge\pi^*\tau\wedge\d t)
        = \pi_*\parens{(\pi^*\d\tau)f\d t + \sum_i\d x_i\pdvof f{x_i}\wedge\pi^*\tau\wedge\d t}\cr
        = \parens{\tau\int f\d t}\d\tau + \parens{\sum_i\int\pdvof f{x_i}\d t}\d x_i\wedge\tau ,
    } $$
    which is equal to $\d\pi_*\omega$.

    Now if $\omega=(\pi^*\tau)f\d t_I$ where $I$ does not contain all $1,\dots,n$.
    Then $\d\pi_*\omega=0$.
    On the other hand
    $$ \d\omega = (\pi^*\d\tau)f\d t_I + (\pi^*\tau)\d f\d t_I. $$
    The first term integrates to zero, but the second term has subterms of the form
    $$ (\pi^*\tau)\pdvof f{t_i}\d t_i\d t_I , $$
    which may integrate to be nonzero when $\d t_i\d t_I=\pm\d t_1\cdots\d t_n$.
    In such a case
    $$ \int_{-\infty}^\infty\pdvof f{t_i}(x,t)\d t_i = f(\dots,+\infty,\dots) - f(\dot,-\infty,\dots) = 0 $$
    as $f$ has compact support.
    And so integrating $\ipdvof f{t_i}\d t$ is zero.
    \qed

\eproof

\bprop[name={projformulas}]

    We have the following projection formulas.
    \benum
        \item Let $\pi\colon E\to M$ be an oriented rank-$n$ vector bundle.
        Let $\tau\in\Omega^\bul(M)$ and $\omega\in\Omega^\bul_\cv(E)$.
        Then
        $$ \pi_*(\pi^*\tau\wedge\omega) = \tau\wedge\pi_*\omega . $$
        \item If further $M$ is an oriented $m$-manifold, $\omega\in\Omega^q_\cv(E)$ and $\tau\in\Omega_c^{m+n-q}(M)$, then with the induced
        orientation on $E$ (by \refmath[lemma]{orientedvboveroriented})
        $$ \int_E(\pi^*\tau)\wedge\omega = \int_M\tau\wedge\pi_*\omega . $$
    \eenum

\eprop

\bproof

    For (1) we may look locally so we may assume $E=M\times\bR^n$.
    Suppose $\omega=\pi^*\eta\cdot f(x,t)\d t_J$ where $J$ does not contain all $1,\dots,n$.
    Then $\pi_*\omega=0$ and
    $$ \pi_*(\pi^*\tau\wedge\omega) = \pi_*(\pi^*(\tau\wedge\eta)f(x,t)\d t_J) = 0 = \tau\wedge\pi_*\omega . $$
    Alternatively suppose $\omega=\pi^*\eta\cdot f(x,t)\d\bar t$.
    Then
    $$ \pi_*(\pi^*\tau\wedge\omega) = (\tau\wedge\eta)\int_{\bR^n}f(x,t)\d\bar t = \tau\wedge\pi_*\omega . $$

    Let $\set{(U_i,\phi_i)}$ be an oriented trivialization for $E$ and $\rho_i$ a partition of unity subordinate to it.
    Then $\omega=\sum_i\rho_i\omega$, so
    $$ \int_E(\pi^*\tau)\wedge\omega = \sum_i\int_{E|_{U_i}}\pi^*\tau\wedge\rho_i\omega $$
    and
    $$ \int_M\tau\wedge\pi_*\omega = \sum_i\int_{U_i}\tau\wedge\pi_*(\rho_i\omega) . $$
    It is thus sufficient to prove this over $E=M\times\bR^n$, which is done as in (1).
    \qed

\eproof

\bthrm[title={Thom Isomorphism}, name={thomiso}]

    Let $\pi\colon E\to M$ be an orientable vector bundle of rank-$n$ over a finite type manifold $M$.
    Then integration along fibers is an isomorphism
    $$ \pi_*\colon H^\bul_\cv(E) \isoto H^{\bul-n}_\dR(M) . $$

\ethrm

\bproof

    Let $U,V$ be open subsets of $M$, so that we can form an exact sequence of complexes
    $$ 0 \to \Omega^\bul_\cv(E|_{U\cup V}) \to \Omega^\bul_\cv(E|_U) \oplus \Omega^\bul_\cv(E|_V) \to \Omega^\bul_\cv(E|_{U\cap V}) \to 0 $$
    as in \refmath{mayerviet}.
    Consider now the diagram
    {\def\diagcolbuf{0pt}\def\diagcolwidth{0pt}
    \diagram{
        \cdots&H^\bul_\cv(E|_{U\cup V})&H^\bul_\cv(E|_U)\oplus H^\bul_\cv(E|_V)&H^\bul_\cv(E|_{U\cap V})&H^{\bul+1}(E|_{U\cup V})&\cdots\cr
        \cdots&H^{\bul-n}(U\cup V)&H^{\bul-n}(U)\oplus H^{\bul-n}(V)&H^{\bul-n}(U\cap V)&H^{\bul+1-n}(U\cup V)&\cdots
    }{
        \diagarrow{from={1,1}, to={1,2}}
        \diagarrow{from={1,2}, to={1,3}}
        \diagarrow{from={1,3}, to={1,4}}
        \diagarrow{from={1,5}, to={1,6}}
        \diagarrow{from={2,1}, to={2,2}}
        \diagarrow{from={2,2}, to={2,3}}
        \diagarrow{from={2,3}, to={2,4}}
        \diagarrow{from={2,5}, to={2,6}}
        \diagarrow{from={1,2}, to={2,2}, text={\pi_*}, x distance=.15cm}
        \diagarrow{from={1,3}, to={2,3}, text={\pi_*}, x distance=.15cm}
        \diagarrow{from={1,4}, to={2,4}, text={\pi_*}, x distance=.15cm}
        \diagarrow{from={1,5}, to={2,5}, text={\pi_*}, x distance=.15cm}
        \diagarrow{from={1,4}, to={1,5}, text={\partial_1}, y distance=.15cm}
        \diagarrow{from={2,4}, to={2,5}, text={\partial_2}, y distance=.15cm}
    }}
    The squares here all commute.
    Commuting of the first two squares follows from linearity of $\pi_*$.
    Recall that
    $$ \pi_*\partial_1\omega = \pi_*\d(\pi^*\rho_V\omega) = \pi_*(\pi^*(\d\rho_V)\wedge\omega + \pi^*\rho_V\d\omega) , $$
    with the projection identity above this is equal to
    $$ = \d\rho_V\wedge\pi_*\omega + \rho_V\wedge\pi_*\d\omega . $$
    And on the other hand
    $$ \partial_2(\pi_*\omega) = \d(\rho_V(\pi_*\omega)) = \d\rho_V\wedge\pi_*\omega + \rho_V\wedge\pi_*\d\omega . $$
    So we have commutativity.

    If $U$ is diffeomorphic to $\bR^n$, then $E|_U$ is trivial by \refmath[proposition]{bundlehomotopy}.
    So we need to show that integration is an isomorphism
    $$ \pi_*\colon H^\bul(M\times\bR^n) \isoto H^{\bul-n}_\dR(M) . $$
    This can be done.
    The proof then reduces to an inductive argument on the size of a good cover of $M$.
    \qed

\eproof

Denote the inverse of $\pi_*$ by
$$ \T\colon H^\bul(M) \to H^{\bul+n}_\cv(E) . $$
The image of $1\in\bR\cong H^0(M)$ in $H^n(E)$ is called the {\emphcolor Thom class} of the oriented bundle $E$.
Notice then that by the projection formula
$$ \pi_*(\pi^*\omega\wedge\T(1)) = \omega\wedge\pi_*\T(1) = \omega . $$
Thus we see that $\T$ can be given by
$$ \T(\bul) = \pi^*(\bul)\wedge\Phi $$
where $\Phi=\T(1)$ is the Thom class of $E$.

\bprop

    The Thom class $\Phi$ on a rank-$n$ oriented bundle $E$ is uniquely characterized as the cohomology class in $H^n_\cv(E)$ which restricts
    to the generator of $H^n_c(F)$ for each fiber $F$.

\eprop

\bproof

    Since $\pi_*\Phi=1$, $\Phi|_F$ is an $n$-form on the fiber with total integral $1$.
    If $\Phi'\in H^n_\cv(E)$ restricts to the generator on each fiber, then
    $$ \pi_*(\pi^*\omega\wedge\Phi') = \omega\wedge\pi_*\Phi' = \omega , $$
    which means that $\pi^*(\bul)\wedge\Phi'$ must be the Thom isomorphism, so $\Phi'=\T(1)=\Phi$.
    \qed

\eproof

\bprop

    If $E,F\to M$ are oriented vector bundles, with projections $\pi_E,\pi_F$, then the Thom class of $E\oplus F$ is
    $$ \Phi(E\oplus F) = \pi_E^*\Phi(E)\wedge\pi_F^*\Phi(F) \in H^{n+m}_\cv(E\oplus F) . $$

\eprop

\bproof

    Let $m$ be the rank of $E$ and $n$ the rank of $F$.
    Then $\pi_E^*\Phi(E)\wedge\pi_F^*\Phi(F)$ is a class of $H^{m+n}_\cv(E\oplus F)$ whose restriction to each fiber is
    the generator of the compact homology of the fiber, via the isomorphism
    $$ H^{m+n}_c(\bR^m\times\bR^n) \cong H^m_c(\bR^m)\otimes H^n_c(\bR^n) $$
    given by the wedge product.
    \qed

\eproof

